\begin{abstract}
Bird’s-Eye-View (BEV) perception serves as a cornerstone for autonomous driving, offering a unified spatial representation that fuses surrounding-view images to enable reasoning for various downstream tasks, such as semantic segmentation, 3D object detection, and motion prediction.
However, most existing BEV perception frameworks adopt an end-to-end training paradigm, where image features are directly transformed into the BEV space and optimized solely through downstream task supervision. 
This formulation treats the entire perception process as a black box, often lacking explicit 3D geometric understanding and interpretability, leading to suboptimal performance.
In this paper, we claim that an explicit 3D representation matters for accurate BEV perception, and we propose \textbf{Splat2BEV}, a Gaussian Splatting-assisted framework for BEV tasks. 
Splat2BEV aims to learn BEV feature representations that are both semantically rich and geometrically precise. 
We first pre-train a Gaussian generator that explicitly reconstructs 3D scenes from multi-view inputs, enabling the generation of geometry-aligned feature representations.
These representations are then projected into the BEV space to serve as inputs for downstream tasks.
Extensive experiments on nuScenes and argoverse dataset demonstrate that Splat2BEV achieves state-of-the-art performance and validate the effectiveness of incorporating explicit 3D reconstruction into BEV perception.
% Extensive experiments on the nuScenes dataset across vehicle, pedestrian, and lane segmentation tasks demonstrate that Splat2BEV achieves state-of-the-art performance 
% and outperforms the existing baselines by a notable 21.4\% improvement on lane segmentation and an average gain of 11.0\% across vehicle, pedestrian and lane segmentation, validating the effectiveness of incorporating explicit 3D reconstruction into BEV perception.
\end{abstract}